\documentclass[aspectratio=169,xcolor=dvipsnames]{beamer}
\usetheme{SimpleDarkBlue}
\setbeamertemplate{caption}[numbered]

\useoutertheme[subsection=false]{smoothbars}
% ^^^ Comment out when deck exceeds ~50 frames (mini-frame dots overflow \paperwidth).
%     Uncomment when the deck is trimmed to a manageable size.
\useinnertheme{circles}
\setbeamertemplate{enumerate items}[default]

% --- Colors ---
\setbeamercolor{section in head/foot}{bg=DarkBlue, fg=white}
\setbeamercolor{section in head/foot shaded}{bg=LightBlue, fg=white}
\setbeamercolor{frametitle}{bg=DarkBlue, fg=white}

\usepackage{url}
\usepackage[authoryear]{natbib}
\bibpunct[; ]{(}{)}{;}{a}{,}{,}
\usepackage{caption}
\usepackage{hyperref}
\usepackage{graphicx}
\usepackage{booktabs}
\usepackage{geometry}

% --- Figure macro: respects both width and height so caption never clips ---
% Usage: \fig{../plots/01_example.pdf}
\newcommand{\fig}[1]{%
  \includegraphics[width=\linewidth,height=0.68\textheight,keepaspectratio]{#1}}

% --- Appendix helper ---
\newcommand{\startappendix}{
  \setcounter{framenumber}{0}
  \setbeamertemplate{footline}{}
  \setbeamertemplate{footline}{
    \hfill\insertframenumber\hspace{2ex}\vspace{2ex}
  }
}
%\journalname{Geophysical Research Letters}

% ----------------------------------------------------------------------------------
%  TITLE PAGE
% ----------------------------------------------------------------------------------
\title[]{ENSO-conditioned evolution of global mean surface temperature}
\author[]{Michael K. Tippett\inst{1} and Emily J. Becker\inst{2}}
\institute[]{\inst{1} Department of Applied Physics and Applied
  Mathematics, Columbia University, New York, New York \and %
  \inst{2} University of Miami Rosenstiel School for Marine, Earth,
  and Atmospheric Science, Miami, Florida}
\date{\today}

\setcounter{secnumdepth}{1}

\AtBeginSection[]{
  \setbeamercolor{background canvas}{bg=blue!50}
  \begin{frame}
    \centering \Huge\textbf{\textcolor{white}{\thesection.\
        \insertsectionhead}}
  \end{frame}
  \setbeamercolor{background canvas}{bg=white}
}
% ----------------------------------------------------------------------------------
%  SLIDES
%  Convention:
%    - Two frames per figure: frame 1 = figure only; frame 2 = figure + bullets
%    - Frame titles: no figure numbers, no citations in captions
%    - Compile from inside tex/ so relative paths (../plots/...) resolve correctly
% ----------------------------------------------------------------------------------
\begin{document}

\begin{frame}
  \titlepage
\end{frame}

% ---- Table of contents (optional) ----
% \begin{frame}{Outline}
%   \tableofcontents
% \end{frame}

\section*{Key points}
\begin{frame}{Key points}
\begin{enumerate}
%\begin{keypoints}
\item A simple regression model quantifies the modulation of monthly
  global mean surface temperature evolution by ENSO

\item Key features captured by the model are the persistence of global
  mean surface temperature and its delayed response to ENSO

\item A December Ni\~no-3.4 anomaly of 1\textdegree C is accompanied
  by global mean temperature increases of 0.03\textdegree C in summer
  and 0.1\textdegree C in late winter
  
%\end{keypoints}
\end{enumerate}
\end{frame}

\section*{Abstract}
\begin{frame}{Abstract}
\begin{enumerate}\scriptsize
\item ENSO forecasts are skillful, and the influence of ENSO on global
  mean surface temperature (GMST) is well established. These facts
  suggest that future GMST behavior can be anticipated based on
  expected ENSO conditions.

\item Here we examined specifically how the June--May trajectory of
  monthly GMST can be anticipated from recent GMST evolution and
  upcoming boreal winter Ni\~no-3.4 values.

\item Principal component analysis and dimension reduction of the GMST
  and Ni\~no-3.4 data led to a simple, interpretable model in which
  the June--May trajectory of monthly GMST is conditioned two
  quantities: the average GMST of the prior 12 months and the upcoming
  December value of Ni\~no-3.4.

\item The simple model improves on previous ones by explicitly
  accounting for serial correlation and seasonality.
  
\item Including ENSO in the simple model improves performance compared
  to baseline models that do not, especially during the period when
  atmospheric bridge mechanisms are active.

\item The simple model allows us to conclude that an expected December
  Ni\~no-3.4 anomaly of 1\textdegree C is associated with a GMST
  increase of about 0.025\textdegree C during boreal summer followed
  by a peak increase of about 0.1\textdegree C in the following
  February.
  
\end{enumerate}
\end{frame}

\section*{Plain Language Summary}
\begin{frame}{Plain Language Summary}
  Global mean surface temperature (GMST) is changing due to
  anthropogenic and natural forcing, as well as natural variability,
  with El Ni\~no-Southern Oscillation (ENSO) being the dominant
  component of natural variability. In this study, we examined how the
  coming June-to-May evolution of global mean temperature can be
  anticipated from two pieces of information available in advance: the
  recent GMST level and the expected state of ENSO over the same
  coming June-to-May period. We developed a simple model that
  associates a December Ni\~no-3.4 anomaly of 1\textdegree C with a
  GMST increase of about 0.025\textdegree C in the preceding boreal
  summer and a peak increase of about 0.1\textdegree C in the
  following February.
\end{frame}


\setcounter{section}{0}
\section{Introduction}

\begin{frame}{Motivation}
  \begin{enumerate}\tiny 
  \item Global mean surface temperature (GMST) is trending upward due
    to human-caused increases in greenhouse gases. Superimposed on
    this upward trend, GMST also varies in response to natural forcing
    (e.g., volcanoes) and internal variability.
  \item ENSO is the dominant component of internal variability on
    seasonal timescales, and its influence on GMST is well established
    \citep[e.g.,][]{Trenberth2002evolution}. During warm ENSO
    conditions (El Ni\~no), GMST tends to be elevated relative to the
    long-term trend, while during cool ENSO conditions (La Ni\~na),
    the opposite tends to occur.
  \item \citet{Lean2008natural} attributed about 0.23\textdegree C of
    GMST warming to the 1997 ``super'' El Ni\~no.
  \item More recently, the World Meteorological Organization noted
    ``The most recent El Ni\~no, in 2023-24, was one of the five
    strongest on record and it played a role in the record global
    temperatures we saw in 2024'' \citep{WMO2026enso}.

  \end{enumerate}

  \begin{enumerate}\tiny
  \item ENSO is predictable, and each month meteorological and climate
    forecast centers around the world issue operational ENSO forecasts
    that extend typically two to four seasons into the future.  ENSO
    forecasts issued after the boreal spring predictability barrier
    are especially skillful \citep[e.g.,][]{Heureux:Chapter}.
  \item In early 2026, forecasts of a strong El Ni\~no have led to
    speculation in the press that it could boost GMST by about
    0.2\textdegree C \citep[e.g.,][]{KeithLucas2026SuperElNino}.
    While this estimate appears plausible, exactly how a given
    forecast of ENSO conditions translates into an expected GMST
    anomaly is less clear, both in terms of timing and amplitude.
  \item For instance, it remains difficult to say, in a quantitative
    and principled manner, what GMST evolution is expected over the
    next 12 months given its recent values and the ENSO conditions
    that are forecast for this winter.
  \item One reason for this difficulty is that previous analyses
    either pooled all months together \citep{Trenberth2002evolution}
    or focused on the ENSO–GMST relationship only at the lag at which
    the correlation is strongest or failed to account explicitly for
    autoregressive behavior \citep{Lean2008natural,Foster2011global}.
  \end{enumerate}
  
  \begin{enumerate}\tiny
  \item Here, we address this gap and present a simple and
    interpretable methodology for anticipating GMST evolution at
    monthly resolution over a June--May 12-month period, given a
    forecast of ENSO conditions over the same period.  We chose this
    calendar period because it matches the ENSO development and decay
    cycle and because ENSO forecasts initialized after boreal spring
    typically have useful skill through the following winter.

  \item We conditioned our GMST analysis on observed ENSO evolution and
    acknowledge imperfect ENSO forecasts as an additional source of
    uncertainty, analogous to the perfect-prog framework in weather
    output statistics \citep{Glahn72,Klein1959objective}.
  \end{enumerate}

  \begin{enumerate}\tiny
  \item Data sources and statistical methods are described in Section
    \ref{sec:data-methods}.  Section \ref{sec:results} introduces our
    data-compression strategies and a sequence of regression models,
    progressing from general formulations to simpler ones that are
    more interpretable and nearly as skillful. We also show why these
    successive model reductions are physically and statistically well
    motivated, and evaluate model performance using hindcasts and
    summary skill measures. Conclusions and discussion are given in
    Section \ref{sec:summary}.
  \end{enumerate}
\end{frame}



\section{Data and Methods}\label{sec:data-methods}

\begin{frame}{Data}
  \begin{enumerate}
  \item Our analysis used monthly data from June 1979 through May
    2025, comprising 46 June--May ENSO years.

  \item Global mean surface temperature (GMST) was computed as the
    area-weighted global mean of gridded surface temperature from NOAA
    GlobalTemp v6.1.0 \citep{Huang2022improvements}. NOAA GlobalTemp
    uses the current WMO climatological standard normals period,
    1991--2020, as the reference period for temperature anomalies. Its
    1850--1900 mean is $-0.77^\circ$C, and this value was subtracted
    to form the preindustrial anomalies shown in the figures.

  \item The Ni\~no-3.4 index was computed from ERSSTv5
    \citep{Huang:ERSST5} as the area-weighted sea surface temperature
    over 5\textdegree S--5\textdegree N, 170\textdegree
    W--120\textdegree W.  Anomalies are shown relative to the
    1991--2020 base period. We made limited use of the relative
    Ni\~no-3.4 index.
  
  \end{enumerate}
\end{frame}

\begin{frame}{Methods}
  \begin{enumerate}\tiny

  \item Principal component analysis (PCA) was applied to the
    $12 \times 46$ Ni\~no-3.4 and GMST data matrices after removal of
    monthly means (row averages). In the data matrices, rows
    correspond to months and columns to years.
    
  \item Double-centered PCA was applied to the same $12 \times 46$
    Ni\~no-3.4 and GMST data matrices after removal of both monthly
    means (row averages) and June--May means (column averages).

  \item We followed the PCA convention in which principal components
    (PCs) have unit variance and empirical orthogonal functions (EOFs)
    carry the physical units \citep{DelsoleTippettBook}.

\item Data were reconstructed as the product of EOFs and PCs, with
  the removed means added back.

\item Coefficient of determination ($R^2$), root-mean-square error
  (RMSE), and mean-squared-error skill score (MSESS) were used as
  summary measures of regression model performance.

  \item Statistical significance of differences in RMSE was assessed
    using the Wilcoxon signed-rank test applied to paired squared
    errors. Because the RMSE values were computed using the same GMST
    data, RMSE values for different models are not independent, and an
    $F$-test for equality of variances is inappropriate
    \citep{TDS:Comparing:2014}. The Wilcoxon signed-rank test
    effectively tests whether the median difference squared errors is
    zero. Statistical significance for RMSE and MSESS is identical
    because MSESS is a monotonic function of RMSE.
    
  \end{enumerate}
\end{frame}

\section{Results}\label{sec:results}

\begin{frame}{Figure 1}
  \begin{minipage}[c]{0.45\linewidth}
    %\fig{/Users/tippett/claude/enso-gmt-prediction/plots/variance_explained.pdf}
    \fig{/Users/tippett/claude/enso-gmt-prediction/plots/eofs_pcs.pdf}
    \fig{/Users/tippett/claude/enso-gmt-prediction/plots/reconstruction_2pc.pdf}
  \end{minipage}\hfill
  \begin{minipage}[c]{0.5\linewidth}
    \captionof{figure}{The leading two EOFs of June--May segments of
      (a) Ni\~no-3.4 and (b) global mean surface temperature (GMST). The
      corresponding leading two PCs of (c) Ni\~no-3.4 and (d) GMST.  The
      standardized December value of Ni\~no-3.4 is also shown in panel
      (c) for comparison with PC1. Two-PC reconstructions of (e)
      Ni\~no-3.4 and (f) GMST with root mean squared error (RMSE) and
      R-squared values in the title.}\label{fig:fig1}
  \end{minipage}
\end{frame}

\begin{frame}{Figure 1}
  \begin{enumerate}\fontsize{6pt}{8pt}\selectfont

  \item A regression model that included all plausible predictors and
    predictands would be high-dimensional: 24 predictors and 12
    predictands. The predictors would be the previous 12 months of
    GMST and the next 12 months of Ni\~no-3.4, while the predictands
    would the next 12 months of GMST. Such a model would be prone to
    overfitting and difficult to interpret. Therefore, as a first
    step, we applied principal component analysis (PCA) to the
    data. PCA provides an unsupervised and systematic dimension
    reduction that is based on explained variance and ignores
    relationships between predictors and predictands.

  \item Previous work has shown that 12-month Ni\~no-3.4 trajectories
    are well approximated by one or two PCs, both in observations and
    in seasonal forecasts \citep{Bunge2009verified,Tippett2019:ENSO}.
    Consistent with these results, we found that approximately 95\% of
    the variability of June--May Ni\~no-3.4 segments is explained by the
    first two PCs (Fig.\ \ref{fig:EOF-variance-explained}a).
    
  \item Moreover, we found that the same is true of
    GMST---approximately 95\% of the variability of a June--May GMST
    trajectory is explained by the first two PCs (Fig.\
    \ref{fig:EOF-variance-explained}b).
  \end{enumerate}
  \begin{enumerate}\fontsize{6pt}{8pt}\selectfont

  \item Although PCA need not yield physically meaningful or readily
    explainable components, the leading PCs of GMST and Niño-3.4
    segments admit relatively clear interpretations.
    
  \item EOF1 of Ni\~no-3.4 peaks in December, reflecting the seasonal
    phase locking of ENSO (Fig.\ \ref{fig:fig1}a).  The correlation of
    PC1 of Ni\~no-3.4 with the December Ni\~no-3.4 index exceeds 0.98
    (compare PC1 of Ni\~no-3.4 with the standardized value of December
    Ni\~no-3.4 in Fig.\ \ref{fig:fig1}c).  EOF2 of Ni\~no-3.4 is a
    dipole in time centered on December, in part because it must be
    orthogonal to EOF1, and in part to account for shifts of the
    Ni\~no-3.4 peak away from December, though the little variance
    explained by PC2 is notable.
    
  \item EOF1 of GMST is nearly uniform across months (Fig.\
    \ref{fig:fig1}b), while PC1 of GMST shows a clear upward trend
    (Fig.\ \ref{fig:fig1}d). This behavior indicates PC1 represents
    the approximately uniform (across calendar months) upward trend of
    GMST.

  \item EOF2 of GMST is nearly uniform and negative June--November
    until it rises in December and peaks in January and February,
    while PC2 shows no visible trends. The asymmetry of its dipole
    structure suggests that PC2 might represent the well-known delayed
    response of GMST to ENSO via the atmospheric bridge
    \citep{Alexander2002atmospheric,Klein1999remote}.
    
  \item Reconstructions of Ni\~no-3.4 and GMST show that 2 EOFs are
    adequate to capture interannual variability ($R^2$ values of about
    95\%; Fig.\ \ref{fig:fig1}e, f), although more unresolved
    subannual variability remains visible in the GMST time series than
    in the Ni\~no-3.4 time series. This difference may reflect the
    fact that much of the GMST variance is associated with its
    long-term trend, which is relatively easily to capture.


    % The two-PC reconstruction of GMST sets an upper bound on the
    % possible accuracy of a two-PC model.

  \end{enumerate}
\end{frame}

\begin{frame}{Figure 2}
  \begin{minipage}[c]{0.55\linewidth}
    %\fig{/Users/tippett/claude/enso-gmt-prediction/plots/variance_explained.pdf}
    \fig{/Users/tippett/claude/enso-gmt-prediction/plots/pc_correlations.pdf}
    \fig{/Users/tippett/claude/enso-gmt-prediction/plots/fit_reduction1.pdf}
  \end{minipage}\hfill
  \begin{minipage}[c]{0.42\linewidth}
    \captionof{figure}{(a) Simultaneous and (b) one-year lag cross
      correlations between the first two PCs of Ni\~no-3.4 and global
      mean surface temperature (GMST).  Panel (c) shows the fit of
      GMST time series by the two-PC model (see text for
      details). Dots show the start of June--May segments. The RMSE
      and R-squared of the four-PC model is shown in the title for
      comparison.}\label{fig:fig2}
  \end{minipage}
\end{frame}

\begin{frame}{Figure 2}
  \begin{enumerate}\scriptsize
  \item We next examined the cross correlations and lag correlations of
    the first two PCs of Ni\~no-3.4 and GMST.
  
  \item The strongest simultaneous correlation is the 0.61 correlation
    between N34-PC1 and GMST-PC2, which supports our hypothesis that
    GMST-EOF2 represents the delayed GMST response to ENSO
    (Fig. \ref{fig:fig2}a). The other simultaneous cross-correlations are
    weak ($R^2 < 4\%$).
  
  \item The strongest lagged relation between PCs is the 0.92
    correlation between GMST-PC1$(t)$ and GMST-PC1$(t+1)$ which
    reflects the upward trend in GMST-PC1 and the resulting strong
    year-to-year persistence (Fig. \ref{fig:fig2}b).  In contrast, the
    lag-1 autocorrelation of GMST-PC2 is small and negative
    ($r=-0.18$), consistent with the interpretation of GMST-PC2 as an
    ENSO response that does not persist from one year to the next.
    
  \item The source of the correlation of N34-PC2$(t)$ with
    N34-PC1$(t+1)$ is less obvious, though it is visually evident in
    the PC time series for both warm and cool ENSO events (Fig.\
    \ref{fig:fig1}c). \textbf{Thoughts?}

  \item The correlation between N34-PC2$(t)$ and GMST-PC2$(t+1)$
    ($r=0.37$) is explained almost entirely by the pathway
    N34-PC2$(t) \rightarrow$ N34-PC1$(t+1) \rightarrow$
    GMST-PC2$(t+1)$. After accounting for N34-PC1$(t+1)$, the partial
    correlation between N34-PC2$(t)$ and GMST-PC2$(t+1)$ is only 0.04.
     
  \item The correlation analysis above suggests that little is lost by
    using two predictors: GMST-PC1$(t)$ and N34-PC1$(t+1)$, instead of
    four: GMST-PC1$(t)$, GMST-PC2$(t)$, N34-PC1$(t+1)$, and
    N34-PC2$(t+1)$. Indeed, the fit obtained with these two predictors
    is nearly identical to that obtained with all four predictors
    (compare Figs.\ \ref{fig:fig2}c and
    \ref{fig:fit_full_recon}). Moreover, in the four-predictor model,
    only the coefficients of GMST-PC1$(t)$ and N34-PC1$(t+1)$ are
    large and statistically significant (Table \ref{tab:full_model}).

  \end{enumerate}
\end{frame}


\begin{frame}{Figure 3}
  \begin{minipage}[c]{0.35\linewidth}
    %\fig{/Users/tippett/claude/enso-gmt-prediction/plots/variance_explained_dc.pdf}
    \fig{/Users/tippett/claude/enso-gmt-prediction/plots/eofs_pcs_dc.pdf}
    %\fig{/Users/tippett/claude/enso-gmt-prediction/plots/enso_response.pdf}
    \fig{/Users/tippett/claude/enso-gmt-prediction/plots/fit_idea2.pdf}
    \fig{/Users/tippett/claude/enso-gmt-prediction/plots/rmse_by_month.pdf}
  \end{minipage}\hfill
  \begin{minipage}[c]{0.42\linewidth}
    \captionof{figure}{Comparison of the standard (a) EOFS and (b) PCs
      of global mean surface temperature (GMST) with double-centered
      ones. Panel (c) shows the fit of GMST time series by the
      simplified model (see text for details). (d) RMSE and (e) MSESS
      of the simplified model compared to persistence
      baselines.}\label{fig:fig3}
  \end{minipage}
\end{frame}

\begin{frame}{Figure 3}
  \begin{enumerate}\fontsize{5pt}{7pt}\selectfont

  \item While the two-PC predictor model is both effective and
    parsimonious, there are further simplifications that can improve
    interpretability.

  \item In particular, the predictors GMST-PC1$(t-1)$ and N34-PC1$(t)$
    are defined through the PCA and therefore depend on both the
    specific data used and the obscure details of the PCA
    calculation. Replacing the PCs with quantities that closely
    approximate them, but whose calculation is more transparent, would
    yield a more explainable model.

  \item First, we replaced N34-PC1 with the December Ni\~no-3.4 value,
    which correlates with N34-PC1 at $r=0.99$ (Fig.\
    \ref{fig:fig1}c). This close correspondence follows from the EOF
    structure: EOF1 peaks in December, while EOF2 is close to zero in
    December.

  \item The resulting model performs essentially as well as the two-PC
    predictor model (Fig.\ \ref{fig:fit_r2_recon}) while being easier
    to interpret. In particular, the dependence of GMST on ENSO can
    now be described without reference to PCs, for example: ``given
    December Ni\~no-3.4 $=x$, the anticipated June--May GMST trajectory
    is \ldots''
  
  \item Second, we replaced GMST-PC1 with the June--May mean GMST
    (which we call the level). The reason for this being that
    GMST-EOF1 is nearly uniform across months. This replacement
    suggests a decomposition of GMST in which the first mode is simply
    the June--May mean, while the second mode captures as much as
    possible of the remaining variance and is expected to resemble
    GMST-PC2.

  \item We implemented this idea by decomposing each June--May GMST
    trajectory into its June--May level together with a PCA of the
    departures from that mean. Because this decomposition removes both
    monthly means and June--May means, we call it double-centered PCA.
    
  \item Double-centered PCA is not variance-optimal so its modes
    explain less variance, but the loss is small: the variance
    explained by the first two modes decreases from 94.8\% to 94.3\%
    (Fig.\ \ref{fig:EOF-variance-explained-dc}). The leading EOF from
    the double-centered decomposition is similar to GMST-EOF2 (Fig.\
    \ref{fig:fig3}a), and the corresponding PC is also similar to
    GMST-PC2, with correlation 0.993 (Fig.\ \ref{fig:fig3}b).

  \item We call the model with these two simplifications (N34-PC1
    replaced by December Ni\~no-3.4 and GMST-PC1 replaced by level), the
    simplified model.
    
  \item The performance of the  simplified model is nearly
    identical to that of the two-PC predictor model (Fig.\
    \ref{fig:fig3}c).

  \item The simplified model improves interpretability by enabling
    statements of the form: ``given recent 12-month average of GMST
    and the expected December Ni\~no-3.4 value, the June--May GMST
    trajectory is expected to be \ldots''
  \end{enumerate}

  \begin{enumerate}\fontsize{5pt}{7pt}\selectfont
  \item A remaining question is how much is gained by conditioning on
    ENSO. Because of the strong upward trend in GMST,
    persistence-based predictions are already a nontrivial
    benchmark. We therefore compared the simplified model with two
    persistence baselines: one based on the preceding June--May average 
    and one based on the preceding May GMST value.

  \item The simplified model has lower RMSE than the June--May-mean
    baseline in every month and lower RMSE than the May baseline from
    August through May (Fig.\ \ref{fig:fig3}d). Its clearest and most
    statistically robust advantage occurs from December through May,
    when the delayed GMST response to ENSO is expected to be
    strongest.

  \item The corresponding MSE skill scores show that the simplified
    model improves substantially on both persistence baselines, with
    the largest gains generally occurring during boreal winter and
    spring (Fig.\ \ref{fig:fig3}e).
  
  \end{enumerate}
\end{frame}

\begin{frame}{Figure 4}
  \begin{minipage}[c]{0.32\linewidth}
    \fig{/Users/tippett/claude/enso-gmt-prediction/plots/hindcast_panels.pdf}
    \fig{/Users/tippett/claude/enso-gmt-prediction/plots/forecast.pdf}
  \end{minipage}\hfill
  \begin{minipage}[c]{0.42\linewidth}
    \captionof{figure}{Observed June--May global mean surface
      temperature (GMST) and simple model fits for (a) 2007-8, (b)
      2009-10, (c) 2010-11, (d) 2015-16, (e) 2020-21, and (f)
      2023-24. Panel (g) shows the implied June--May GMST response per
      $+1^\circ$C of December Ni\~no-3.4 and relative
      Ni\~no-3.4. Panel (h) shows a 2026-27 illustration for Dec 2026
      Ni\~no-3.4 of 2.0\textdegree C.}\label{fig:fig4}
  \end{minipage}
\end{frame}

\begin{frame}{Figure 4}  
  \begin{enumerate}\tiny
  \item Summary performance metrics are informative, but they do not
    show how the simplified model compares to observations in
    individual years.  To address this point, Figs.\
    \ref{fig:fig4}a--f shows representative hindcasts for several
    recent years.

\item The prediction intervals are fairly wide, but in most cases
  they provide a reasonable characterization of uncertainty.

\item In the cool ENSO years shown (2007, 2010, and 2020), both
  the simplified model and the observed GMST generally lie below
  the level of the previous 12 months, with a pronounced dip
  around January and a partial rebound by May. The amplitude of
  the simplified model is smaller than that of the observed GMST,
  as expected from a regression model; in particular, the
  relevant coefficient of determination is only $R^2=0.40$
  (Table \ref{tab:model-sm}).

\item In the warm ENSO cases, the behavior is broadly opposite,
  with elevated GMST during boreal winter and early spring.
  An important exception is 2023--24, when the sharp warming in
  September is not consistent with the expected timing of the
  delayed ENSO influence and lies outside the prediction
  intervals.

\item An attractive feature of the simplified model is that, once the
  recent GMST level is specified, the month-to-month shape of the
  predicted GMST response seasonal pattern is fixed and only its
  amplitude changes with the December Ni\~no-3.4 value.

\item This structure makes it straightforward to describe the GMST
  response to a specified ENSO condition. Figure \ref{fig:fig4}g
  shows the implied June--May GMST response per
  $+1^\circ$C of December Ni\~no-3.4. The response is about
  $0.025^\circ$C during boreal summer, rises sharply beginning in
  December, and peaks at slightly more than $0.10^\circ$C in
  February.

\item The same analysis can be carried out using relative Ni\~no-3.4
  \citep{LHeureux2024relative}. The implied GMST response is similar,
  although slightly weaker in amplitude, which might be due to the
  fact that the variance of December relative Ni\~no-3.4 is about 8\%
  greater than that of the traditional Ni\~no-3.4 index over this
  period. Thus, a 1\textdegree C anomaly of relative Ni\~no-3.4 is
  ``smaller'' than a 1\textdegree C anomaly of traditional
  Ni\~no-3.4. Note the correlations are abou the same
  (\ref{fig:dec-n34-corr}).
  
\item We also show an illustrative 2026--27 GMST trajectory based
  on the currently observed June 2025--March 2026 mean of
  1.27$^\circ$C above 1850--1900 together with an assumed
  December 2026 Ni\~no-3.4 value of $+2.0^\circ$C.

\item Under these conditions, the predicted trajectory peaks at about
  1.5$^\circ$C above preindustrial in February 2027, close to the
  expected seasonal peak of ENSO influence.
\end{enumerate}
\end{frame}


\section{Summary and discussion}\label{sec:summary}

\begin{frame}[allowframebreaks]{Conclusions}
  
  \begin{enumerate}\fontsize{6pt}{8pt}\selectfont
  \item \textbf{Summary}

  \item Here we developed a simple and interpretable regression model
    that quantifies the well-known relationship between ENSO and
    global mean surface temperature (GMST). The model is predictive in
    that it allows the June--May evolution of monthly GMST to be
    anticipated based on information that is available near the
    beginning of June.

  \item We began with a high-dimensional formulation and applied a
    sequence of systematic, physically motivated simplifications to
    obtain a model with two predictors. The resulting
    predictors---recent GMST level and expected December Ni\~no-3.4
    value---serve to model the persistence of GMST and its delayed
    response to ENSO.

  \item These reductions were informed by principal component analysis
    (PCA) of 12-month June--May trajectories of GMST and Ni\~no-3.4.

  \item For both GMST and Ni\~no-3.4, the first two PCs explain more
    than 95\% of the variance of June--May trajectories.

  \item Explaining variance is notably easier for GMST than for
    Ni\~no-3.4 because a substantial fraction of GMST variance is
    explained by a long-term trend.  For Ni\~no-3.4, by contrast, the
    result is more striking, since a common view is that each ENSO
    event has its own distinct temporal evolution.
  
  \item The structure of the leading EOFs of the data suggested
    further simplifications to improve interpretability. Since the
    leading EOF of GMST is nearly uniform across months, it was
    replaced by the June--May average of GMST.

  \item The leading Ni\~no-3.4 EOF represents the typical temporal
    evolution of ENSO and peaks in December, so its PC was replaced by
    the December value of Ni\~no-3.4.
    
  \item These changes simplified the predictors with little loss of
    accuracy.
    
  \item According to the model, an expected December Ni\~no-3.4
    anomaly of 1\textdegree C is associated with a GMST increase of
    about 0.025\textdegree C in boreal summer and a peak increase of
    about 0.1\textdegree C in the following February.

  \item Although the influence of ENSO on GMST is modest, including
    ENSO in the model increased its accuracy compared to
    persistence-only baselines.
  
    % \item The simple model quantifies a familiar picture GMST
    %   evolution: long-term warming sets the background, while ENSO
    %   events drives departures from that background a few months
    %   after the ENSO peak in boreal winter.

\end{enumerate}

\begin{enumerate}\fontsize{6pt}{8pt}\selectfont
\item \textbf{How this work is different.}

\item Regression-based analyses of ENSO and global temperature have
  been carried out previously, but they were designed for different
  purposes and used different model structures.

\item In particular, earlier studies were often aimed at removing
  ENSO-related variability from GMST to better isolate trends and
  responses to external forcing.

\item A key difference between the model here and earlier approaches
  is that earlier models were generally not explicitly autoregressive
  in GMST. In some cases ENSO was the only predictor
  \citep{Angell2000tropospheric,Trenberth2002evolution}.  In others,
  linear trends in time or terms representing external forcing were
  included, but lagged GMST was not
  \citep{Foster2011global,Lean2008natural}.

\item Earlier analyses also tended not to be explicitly seasonal.  For
  example, previous work often pooled months together and concluded
  that the ENSO--GMST relation peaked at a lag of about three months
  \citep{Trenberth2002evolution}. Here, the influence of ENSO on GMST
  is more precisely described not as ``occurring at a lag of a few
  months'' but as ``occurring at a lag of a few months after the
  December peak,'' that is, a seasonally phased response, consistent
  with the atmospheric-bridge mechanisms
  \citep{Alexander2002atmospheric,Klein1999remote}.

\end{enumerate}


\begin{enumerate}\fontsize{6pt}{8pt}\selectfont

\item \textbf{GCMs}

\item A natural question is GCMs fit into all this. They represent the
  coupled climate system and provide a physically self-consistent
  description of ENSO and GMST evolution.

\item Climate projections, such as those used in IPCC assessments, are
  not expected to represent the observed year-to-year influence of
  ENSO on GMST because they are not initialized with observed
  conditions. Consequently their representations of ENSO are neither
  in phase with reality nor across model and members.

  
\item Decadal predictions from the WMO Lead Centre for Annual to
  Decadal Climate Prediction start from the conditions on 1 November
  or 1 January initial conditions, which means they can only take
  advantage of the ENSO signal for a few months until they run into
  the spring predictability barrier
  \citep{Hermanson2022wmo,WMO2025GADCU}.

  
\item Operational annual forecasts of calendar-year global mean
  temperature have the same problem.

\item UK Met Office ``DePreSys3 global mean temperature predictions,
  initialized every 1st November'' \citep{Dunstone2024will}.
  
\item Similar issues with UK Met Office statistical forecasts that use
  a October–November ENSO index which doesn't get them far
  \citep{Folland2025review}. Also the ENSO index is based on the high
  frequency SST EOF1 which is obscure af.

\item Initialized seasonal climate models provide a physically
  self-consistent way to predict GMST, because they
  evolve the coupled climate system forward rather than fitting global
  mean temperature statistically after the fact.

\item For the Copernicus Climate Change Service (C3S), the standard
  lead time for seasonal forecast products is six months.
  
\item However, \citet{Tippett2024trends} showed that current systems
  still have important shortcomings for this problem, even though they
  do exhibit skill beyond a simple trend baseline.

\item In particular, much of that skill is tied to ENSO, but the
  forecast quality depends strongly on season, forecast uncertainty is
  sensitive to initialization, individual models tend to be
  underdispersed at short lead, and the multimodel ensemble tends to
  be overdispersed.

\item Initialized seasonal climate forecast systems do predict GMST,
  but systematic assessments of that capability have been relatively
  uncommon. In particular, recent work with the North American
  Multi-Model Ensemble (NMME) found that initialized seasonal
  forecasts of monthly global mean temperature do have skill beyond
  that of the trend and that much of that additional skill is
  attributable to ENSO. The same study also found important
  deficiencies, including seasonal dependence of skill (SPRING
  predictability barrier!!!), underdispersion in single-model
  ensembles, and overdispersion in the multi-model ensemble

\item There is a useful analogy here with the older two-tier seasonal
  prediction strategy used widely in the 1990s and 2000s, in which SST
  was predicted first and the atmospheric response was then computed
  separately with an AGCM, in part because coupled-model biases were
  often too large to rely on a fully coupled forecast
  alone. 

\item Our approach is similar in spirit. Rather than relying on a full
  coupled prediction system to get every part of the problem right at
  once, we isolate the most predictable large-scale components and use
  them in a simpler predictive framework.

\item In particular, we use the slowly varying GMST background state
  together with the seasonally dependent ENSO influence to anticipate
  the June--May evolution of monthly GMST.

\end{enumerate}

\begin{enumerate}\fontsize{6pt}{8pt}\selectfont
\item \textbf{Future work}
\item Explore alternative representations of ENSO, including indices
  that distinguish event amplitude, longitude, or temporal evolution,
  to determine whether they better capture the seasonally varying
  GMST response.

\item Extend the framework to initialized ENSO forecasts, so that the
  December Ni\~no-3.4 predictor is itself predicted operationally rather
  than prescribed.

\item Compare the statistical framework more directly with initialized
  seasonal climate model forecasts of GMST, with
  particular attention to reliability, dispersion, and sensitivity to
  initialization.

\item Investigate whether the same low-dimensional trajectory approach
  can be applied to regional or zonal mean temperatures, where the
  forced signal is weaker and the seasonal ENSO response may be more
  structured.

\end{enumerate}
\end{frame}

%%%%%%%%%%%%%%%%%%%%%%%%%%%%%%%%%%%%%%%%%%%%%%%
%
% DATA SECTION and ACKNOWLEDGMENTS
%
%%%%%%%%%%%%%%%%%%%%%%%%%%%%%%%%%%%%%%%%%%%%%%%

%% AGU_OPENRESEARCH_BEGIN
\begin{frame}{Open Research}
  \begin{enumerate}
\item NOAAGlobalTemp version 6.1 data are from
  \url{https://www.ncei.noaa.gov/data/noaa-global-surface-temperature/v6/access/gridded/NOAAGlobalTemp_v6.1.0_gridded_s185001_e202602_c20260308T114550.nc}.
\item NOAA Extended Reconstructed SST V5 (ERSSTv5) data are provided
  by the NOAA PSL, Boulder, Colorado, USA, from their website
  \url{https://psl.noaa.gov}.
\item Monthly values of the relative Ni\~no-3.4 index are available
  from
  \url{https://www.cpc.ncep.noaa.gov/data/indices/Rnino34.ascii.txt}.
\end{enumerate}
\end{frame}
%% AGU_OPENRESEARCH_END

%% AGU_COI_BEGIN
\begin{frame}{Conflict of Interest}
  The authors declare no conflicts of interest.
\end{frame}
%% AGU_COI_END


%% AGU_ACKS_BEGIN
\begin{frame}{Acknowledgments}
  Supported by NSF AGS-0000000.
\end{frame}
%% AGU_ACKS_END
     
\begin{frame}[allowframebreaks]{References}
  \scriptsize
  \bibliographystyle{ametsocV6}
  \bibliography{all}
\end{frame}

\clearpage

\section{Supplemental/not shown}

\renewcommand\thefigure{S\arabic{figure}}
\setcounter{figure}{0}

\renewcommand\thetable{S\arabic{table}}
\setcounter{table}{0}

% ---- Variance explained ----
\begin{frame}{Variance explained by EOF mode}
  \begin{minipage}[c]{0.55\linewidth}
    \fig{/Users/tippett/claude/enso-gmt-prediction/plots/variance_explained.pdf}
  \end{minipage}\hfill
  \begin{minipage}[c]{0.42\linewidth}
    \captionof{figure}{Variance per mode and cumulative explained
      variance of PCA applied to (a, c) Ni\~no-3.4 and (b, d) global
      mean temperature (GMST). The dashed lines in panels a, b indicate
      the 5\% level and in panels c, d the 90\%
      level.}\label{fig:EOF-variance-explained}
  \end{minipage}
\end{frame}

% ---- Full model fit ----
\begin{frame}{Full conditional regression: benchmark}
  \begin{minipage}[c]{0.55\linewidth}
    \fig{/Users/tippett/claude/enso-gmt-prediction/plots/fit_full_model.pdf}
  \end{minipage}\hfill
  \begin{minipage}[c]{0.42\linewidth}
    \captionof{figure}{Reconstruction fit of the full two-PC
      model.}\label{fig:fit_full_recon}
  \end{minipage}
\end{frame}

\begin{frame}{December Ni\~no-3.4 as predictor instead of N34-PC1}
  \begin{minipage}[c]{0.55\linewidth}
    \fig{/Users/tippett/claude/enso-gmt-prediction/plots/fit_reduction2.pdf}
  \end{minipage}\hfill
  \begin{minipage}[c]{0.42\linewidth}
  \captionof{figure}{Reconstruction fit of reduction 2.}\label{fig:fit_r2_recon}
  \end{minipage}
\end{frame}


% ---- Idea 2: double-centered decomposition ----
\begin{frame}{GMST variance partition: level vs shape}
  \begin{minipage}[c]{0.55\linewidth}
    \fig{/Users/tippett/claude/enso-gmt-prediction/plots/variance_explained_dc.pdf}
  \end{minipage}\hfill
  \begin{minipage}[c]{0.42\linewidth}
    \captionof{figure}{Variance per mode and cumulative
      variance for double-centered PCA.}\label{fig:EOF-variance-explained-dc}
  \end{minipage}
\end{frame}

% ---- Correlation of Dec Nino-3.4 with DC residual ----
\begin{frame}{Dec Ni\~no-3.4 correlation with DC GMST residual by month}
  \begin{minipage}[c]{0.55\linewidth}
    \fig{/Users/tippett/claude/enso-gmt-prediction/plots/corr_dec_n34_by_month.pdf}
  \end{minipage}\hfill
  \begin{minipage}[c]{0.42\linewidth}
    \captionof{figure}{Dec Ni\~no-3.4 correlation with DC GMST
      residual by month}\label{fig:corr_res}
    \begin{itemize}\tiny
    \item Figure \ref{fig:corr_res}: Expected: positive correlations
      peaking near December (ENSO peak season) --- confirmed.
      Unexpected: Jun--Aug correlations are equally strong but
      \emph{negative}.
    \item The small Jun--Aug regression coefficients (see left panel
      of previous figure) are not sampling noise --- they represent
      genuine signal.
    \item Hypothesis: ENSO events tend to alternate in sign year to
      year.  When Dec Ni\~no-3.4 is warm (El Ni\~no), the preceding
      Jun--Aug was more likely cool (residual La Ni\~na from the prior
      year), producing negative correlation.
    \item Implication: the DC residual in early months carries memory
      of the \emph{prior} ENSO event, not captured by Dec Ni\~no-3.4
      alone.
    \end{itemize}
  \end{minipage}
\end{frame}


\begin{frame}{Idea~2 model: predictor--predictand correlations}
  \begin{minipage}[c]{0.55\linewidth}
    \fig{/Users/tippett/claude/enso-gmt-prediction/plots/pc_correlations_dc_pred.pdf}
  \end{minipage}\hfill
  \begin{minipage}[c]{0.42\linewidth}
    \captionof{figure}{Cross correlations for idea 2.}\label{fig:idea-2-cross-corr}
  \end{minipage}
\end{frame}

% ---- Cross-correlations (double-centered model) ----
\begin{frame}{Cross-correlations: level and shape predictands}
  \begin{minipage}[c]{0.55\linewidth}
    \fig{/Users/tippett/claude/enso-gmt-prediction/plots/pc_correlations_dc.pdf}
  \end{minipage}\hfill
  \begin{minipage}[c]{0.42\linewidth}
    \captionof{figure}{Cross-correlations: level and shape predictands traditional and relative Nino-3.4}\label{fig:dec-n34-corr}
  \end{minipage}
\end{frame}


\begin{frame}{Four-PC model: regression table}
  \centering
\begin{table}
  \input{/Users/tippett/claude/enso-gmt-prediction/tex/regression_table_full.tex}
\caption{Regression table for the four-PC model}\label{tab:full_model}
\end{table}
\end{frame}

\begin{frame}{Two-PC model: regression table}
\begin{table}
  \centering
  \input{/Users/tippett/claude/enso-gmt-prediction/tex/regression_table_r1.tex}
  \caption{Regression table for the two-PC model}\label{tab:model-r1}
  \end{table}
\end{frame}

\begin{frame}{N34-PC1 replaced by Dec Ni\~no-3.4: regression table}
\begin{table}
  \centering
  \input{/Users/tippett/claude/enso-gmt-prediction/tex/regression_table_r2.tex}
  \caption{Regression table for N34-PC1 replaced by Dec Ni\~no-3.4 in the two-PC model.}
\end{table}
\end{frame}

\begin{frame}{Simplified model: regression table}
\begin{table}
  \centering
  \input{/Users/tippett/claude/enso-gmt-prediction/tex/regression_table_dc2.tex}
  \caption{Regression table for the simplified model.}
  \label{tab:model-sm}
\end{table}
\end{frame}

\begin{frame}{Simplified model: regression table}
\begin{table}
  \centering
  \input{/Users/tippett/claude/enso-gmt-prediction/tex/regression_table_dc2_roni.tex}
  \caption{Regression table for the simplified model but with RONI.}
  \label{tab:model-sm-roni}
\end{table}
\end{frame}

\end{document}
